% Chapter introduction: why cohomology arose, and what it answers.
% STATUS: written 2026-11-06 (English gate clean; sources read in full as recorded below).
Chapter~\ref{ch:origins} placed derivations inside a cohomology theory. In 1945 Hochschild
set out to connect the structure of an algebra to the vanishing of its cohomology
groups~\cite[p.~58]{Hochschild1945}. His $1$-cocycles are derivations, his
$1$-coboundaries are inner derivations, and the separable algebras are exactly those
on which every derivation is inner. This chapter follows the theory from that point.
\hbfigure{19-cohomology/hochschild-low}{The Hochschild complex in low degrees: derivations are the one-cocycles, inner derivations the one-coboundaries.}{fig:coh:low}

Section~\ref{sec:cohomology:first-cohomology} gives the first group $H^1$ its modern
form, derivations modulo inner derivations, and recalls why it controls the rest.
Section~\ref{sec:cohomology:computations} computes the theory for tensor algebras,
commutative algebras and matrix rings. It also reports the theory built for an
algebra together with a derivation. This covers the generalized derivation
extensions of $n$-Lie algebras
\cite[Proposition~2.13, pp.~25--26]{AtesliEsenSutlu2021}, the
$3$-LieDer pairs of Xu and Liu \cite[p.~1]{XuLiu2021ThreeLieDer}, and the
non-abelian cohomology of Nijenhuis Lie algebras
\cite[Theorem~3.4, p.~6]{Das2025Nijenhuis}. Section~\ref{sec:cohomology:algebraic-topology} reads the same
rules as Hochschild and cyclic homology of a simplicial object, and compares
them with de Rham cohomology.
Section~\ref{sec:cohomology:gerstenhaber} turns to
deformations, where $H^2$ and $H^3$ govern how an algebra can be changed.
Section~\ref{sec:cohomology:noncommutative-geometry} then keeps a differential
calculus when the functions no longer commute.
Section~\ref{sec:cohomology:bounded-cohomology} passes to operator algebras, where only
bounded cochains are admitted and the theory becomes trivial.
Section~\ref{sec:cohomology:group-cohomology} closes the chapter with the parallel theory
of groups, where crossed homomorphisms play the role of derivations.
